\section{Covariance matrices}
\label{sec:covariance}
For $n\geq0$, use coordinates
$(q_1,\ldots,q_n,p_1,\ldots,p_n)$ on $\R^{2n}$ and define
\[
 \Omega_n=\begin{pmatrix}0&I_n\\-I_n&0\end{pmatrix}.
\]
For a complex matrix $H$, the notation $H\succeq0$ means that $H=H^*$ and
$z^*Hz\geq0$ for every complex column vector $z$. Here $H^*$ is the conjugate
transpose. Real matrices are regarded as complex matrices when used in this
inequality.

\begin{definition}[Admissible covariance]\label{def:admissible}
A real $2n\times2n$ matrix $\gamma$ is called \emph{admissible} if
\begin{equation}\label{eq:admissibility}
 \gamma^T=\gamma,\qquad \gamma-i\Omega_n\succeq0.
\end{equation}
\end{definition}
These coordinates are dimensionless. For one mode the condition concerns the
Hermitian matrix
\[
 \begin{pmatrix}a&c-i\\c+i&b\end{pmatrix},
 \qquad \gamma=\begin{pmatrix}a&c\\c&b\end{pmatrix}.
\]
The imaginary part records the canonical pairing between position and momentum.
The empty matrix is admissible when $n=0$.

\subsection{Changes of coordinates}
A real matrix $S$ is symplectic in these coordinates when
$S\Omega_nS^T=\Omega_n$. Transforming a covariance in both indices gives
$\gamma\mapsto S\gamma S^T$. This is a congruence transformation. Similarity
has the form $A\mapsto SAS^{-1}$; the two formulas agree for orthogonal $S$,
but symplecticity alone does not give $S^T=S^{-1}$.

\begin{theorem}\label{thm:congruence}
Let $n\geq0$ and let $S,\gamma$ be real $2n\times2n$ matrices. If $\gamma$
is admissible and $S\Omega_nS^T=\Omega_n$, then $S\gamma S^T$ is admissible.
\end{theorem}
\begin{proof}
Symmetry follows from $(S\gamma S^T)^T=S\gamma^TS^T=S\gamma S^T$.
Put $H=\gamma-i\Omega_n$. Since $S$ is real,
\[
 S\gamma S^T-i\Omega_n=S(\gamma-i\Omega_n)S^T=SHS^*.
\]
This matrix is Hermitian, and for every $z\in\C^{2n}$,
\[
 z^*SHS^*z=(S^*z)^*H(S^*z)\geq0.
\]
The same identities hold for empty matrices.
\end{proof}

\subsection{Subsystems and independent sums}
To select a subsystem, a position coordinate and its matching momentum
coordinate must be selected together. A \emph{one-to-one map}, also called an
\emph{injective map} or an \emph{injection}, is a map $e$ such that
$e(i)=e(j)$ implies $i=j$. Thus a one-to-one map
$e:\{1,\ldots,m\}\to\{1,\ldots,n\}$ selects distinct modes without
repetitions, in the order
\[
 (q_{e(1)},\ldots,q_{e(m)},p_{e(1)},\ldots,p_{e(m)}).
\]
Write $E(q_j)=q_{e(j)}$ and $E(p_j)=p_{e(j)}$. The selected covariance has
entries $\gamma'_{ab}=\gamma_{E(a),E(b)}$.

\begin{theorem}\label{thm:subsystem}
For all $m,n\geq0$, every such selection from an admissible $n$-mode covariance
is an admissible $m$-mode covariance. The empty selection is included.
\end{theorem}
\begin{proof}
Symmetry is inherited by restriction. The $qq$ and $pp$ blocks of $\Omega_n$
vanish, and its $qp$ and $pq$ entries are $\delta_{ij}$ and $-\delta_{ij}$.
Injectivity gives $\delta_{e(i),e(j)}=\delta_{ij}$, so restriction along $E$
sends $\Omega_n$ to $\Omega_m$.

Thus $H'=\gamma'-i\Omega_m$ is the selected principal submatrix of
$H=\gamma-i\Omega_n$. It is Hermitian. Extending a vector $z\in\C^{2m}$
by zero on the unselected coordinates gives $\widetilde z\in\C^{2n}$ and
\[
 z^*H'z=\widetilde z^{\,*}H\widetilde z\geq0.
\]
\end{proof}

Conversely, independent covariance blocks can be combined without cross terms.
\begin{definition}[Independent sum]\label{def:independent-sum}
For arbitrary real matrices $A$ and $B$ of sizes $2n\times2n$ and
$2m\times2m$, write their position/momentum blocks as $A_{qq},A_{qp},A_{pq},A_{pp}$
and similarly for $B$. Define
\begin{equation}\label{eq:independent-sum}
 A\boxplus B=
 \begin{pmatrix}
 A_{qq}&0&A_{qp}&0\\
 0&B_{qq}&0&B_{qp}\\
 A_{pq}&0&A_{pp}&0\\
 0&B_{pq}&0&B_{pp}
 \end{pmatrix}
\end{equation}
in the order $(q_A,q_B,p_A,p_B)$. This is $\diag(A,B)$ reordered from
$(q_A,p_A,q_B,p_B)$; the definition makes no symmetry assumption.
\end{definition}

\begin{theorem}\label{thm:sum}
For all $n,m\geq0$ and all real matrices $A,B$ of the stated sizes,
$A\boxplus B$ is admissible if and only if both $A$ and $B$ are admissible.
\end{theorem}
\begin{proof}
The same coordinate permutation sends $\diag(\Omega_n,\Omega_m)$ to
$\Omega_{n+m}$. Hence the uncertainty matrix of $A\boxplus B$ is the
corresponding permutation of $\diag(A-i\Omega_n,B-i\Omega_m)$.
When both blocks are admissible, symmetry is preserved and the quadratic
form is the sum of two nonnegative quadratic forms. Hermiticity is inherited
from the blocks. Conversely, restriction to either collection of matching
position/momentum coordinates recovers that block and its uncertainty matrix.
Theorem~\ref{thm:subsystem} applies. Empty blocks cause no exception.
\end{proof}

\subsection{One-mode diagonal matrices}
The diagonal case gives explicit examples and fixes the scale of the uncertainty
bound. Put $D(a,b)=\diag(a,b)$ in the order $(q,p)$.
\begin{proposition}\label{prop:diagonal}
For all $a,b\in\R$, the matrix $D(a,b)$ is admissible if and only if
$a>0$ and $ab\geq1$.
\end{proposition}
\begin{proof}
The uncertainty matrix is
$H=\left(\begin{smallmatrix}a&-i\\i&b\end{smallmatrix}\right)$.
If $H\succeq0$, the coordinate vectors give $a,b\geq0$. If $a=0$,
the vector $(b+1,-i)^T$ gives the negative value $-b-2$; hence $a>0$.
The vector $(i,a)^T$ then gives $0\leq a(ab-1)$, and therefore $ab\geq1$.

For the converse, suppose $a>0$ and $ab\geq1$. Write $z=x+iy$, $w=u+iv$.
The quadratic form is real and satisfies
\begin{align*}
 Q&=(\overline z,\overline w)H(z,w)^T
   =a(x^2+y^2)+b(u^2+v^2)+2(xv-yu),\\
 aQ&=(ax+v)^2+(ay-u)^2+(ab-1)(u^2+v^2)\geq0.
\end{align*}
Thus $Q\geq0$; Hermiticity of $H$ is immediate.
\end{proof}

In particular, $I_2=D(1,1)$ is admissible. More generally,
$D(t,t^{-1})$ is admissible for every $t>0$. Its two variances in our covariance
convention are $t/2$ and $1/(2t)$; their product is $1/4$.
By contrast, $cI_2$ with $0\leq c<1$ is a positive semidefinite real matrix
but is not admissible, since $c^2<1$.
Theorem~\ref{thm:sum} combines admissible one-mode blocks into larger matrices.
We now construct wave functions whose covariances give the one-mode equality case.
